% GENERATED FILE. Edit canonical lessons, then run npm run generate:books.
% canonical-source-hash: fnv1a64:e3dc15ceb2f77012
% canonical-lessons: RU-C331-testo, RU-C331-mesit, RU-C331-raskatyvat, RU-C331-lepit, RU-C331-nachinka

\chapter{Dough}
\label{ch:ru-dough}
\hlchaptermodality{331}

\begin{chapteropening}
\textbf{By the end of this chapter:} \emph{I can describe making dough: dough, kneading, rolling out, shaping and a filling.}

Complete the last lesson of chapter 331: I can describe making dough: dough, kneading, rolling out, shaping and a filling.
\end{chapteropening}

\section[\texorpdfstring{\ru{тесто} (tésto)}{тесто (tésto)}]{\ru{тесто} (tésto) — dough}

\label{lesson:RU-C331-testo}



\begin{quote}
Before the new one: say the Russian for to stir, then the Russian for a ladle.
\end{quote}

\subsection*{You'll want to know: \ru{тесто}}
\textbf{\ru{тесто}} (\emph{tésto}) — \textquotedblleft{}dough\textquotedblright{}.

\begin{quote}
\ru{Тесто} \ru{должно} \ru{постоять} \ru{час}.
\end{quote}

\begin{quote}
\emph{Tésto dolzhnó postoyát chas.}
\end{quote}

\begin{quote}
The dough has to rest for an hour.
\end{quote}

Neuter.

\begin{cousinweb}
\emph{\ru{тесто}} is inherited through Old East Slavic \emph{těsto} from Proto-Slavic \emph{těsto} \textquotedblleft{}dough\textquotedblright{} (EtymDB).
\end{cousinweb}

\begin{grammarlens}[title={What you've built}]
One more word for telling, explaining and describing at B1.
\end{grammarlens}

\subsection*{Your turn}
\begin{itemize}
\raggedright
  \item \emph{Say it:} \emph{tésto}
  \item \emph{Say it:} \emph{tésto}, once more
  \item \emph{Say it:} a sentence of your own with \emph{tésto}
  \item \emph{Recall:} say the Russian for to stir, then the Russian for a ladle, then say \emph{tésto} again
\end{itemize}

\begin{tcolorbox}[breakable,colback=teal!4,colframe=teal!35!black,title={Before you move on}]
What does \ru{тесто} mean? (\textquotedblleft{}Dough\textquotedblright{}.) Say it once more.
\end{tcolorbox}
\section[\texorpdfstring{\ru{месить} (mesít)}{месить (mesít)}]{\ru{месить} (mesít) — to knead}

\label{lesson:RU-C331-mesit}



\begin{quote}
Before the new one: say the Russian for a ladle, then the Russian for dough.
\end{quote}

\subsection*{You'll want to know: \ru{месить}}
\textbf{\ru{месить}} (\emph{mesít}) — \textquotedblleft{}to knead\textquotedblright{}.

\begin{quote}
\ru{Бабушка} \ru{месит} \ru{тесто}.
\end{quote}

\begin{quote}
\emph{Bábushka mésit tésto.}
\end{quote}

\begin{quote}
Grandma is kneading the dough.
\end{quote}

Imperfective; present \emph{\ru{мешу}}, \emph{\ru{месишь}}. Perfective \emph{\ru{замесить}} (\emph{zamesít}).

\begin{cousinweb}
\emph{\ru{месить}} is inherited from Proto-Slavic \emph{měsiti}, from PIE \emph{mey\'{k}-} \textquotedblleft{}to mix\textquotedblright{} (EtymDB).
\end{cousinweb}

\begin{grammarlens}[title={What you've built}]
One more word for telling, explaining and describing at B1.
\end{grammarlens}

\subsection*{Your turn}
\begin{itemize}
\raggedright
  \item \emph{Say it:} \emph{mesít}
  \item \emph{Say it:} \emph{mesít}, once more
  \item \emph{Say it:} a sentence of your own with \emph{mesít}
  \item \emph{Recall:} say the Russian for a ladle, then the Russian for dough, then say \emph{mesít} again
\end{itemize}

\begin{tcolorbox}[breakable,colback=teal!4,colframe=teal!35!black,title={Before you move on}]
What does \ru{месить} mean? (\textquotedblleft{}To knead\textquotedblright{}.) Say it once more.
\end{tcolorbox}
\section[\texorpdfstring{\ru{раскатывать} (raskátyvat)}{раскатывать (raskátyvat)}]{\ru{раскатывать} (raskátyvat) — to roll out}

\label{lesson:RU-C331-raskatyvat}



\begin{quote}
Before the new one: say the Russian for dough, then the Russian for to knead.
\end{quote}

\subsection*{You'll want to know: \ru{раскатывать}}
\textbf{\ru{раскатывать}} (\emph{raskátyvat}) — \textquotedblleft{}to roll out\textquotedblright{}.

\begin{quote}
\ru{Он} \ru{раскатал} \ru{тесто} \ru{на} \ru{столе}.
\end{quote}

\begin{quote}
\emph{On raskatál tésto na stolé.}
\end{quote}

\begin{quote}
He rolled out the dough on a table.
\end{quote}

Imperfective; perfective \emph{\ru{раскатать}} (\emph{raskatát}).

\begin{cousinweb}
\emph{\ru{раскатывать}} is built on \emph{\ru{рас}-} \textquotedblleft{}apart\textquotedblright{} and \emph{\ru{катать}} \textquotedblleft{}to roll\textquotedblright{}, an inherited Slavic verb (EtymDB lists \emph{\ru{катать}} \textquotedblleft{}to roll\textquotedblright{} with no parent).
\end{cousinweb}

\begin{grammarlens}[title={What you've built}]
One more word for telling, explaining and describing at B1.
\end{grammarlens}

\subsection*{Your turn}
\begin{itemize}
\raggedright
  \item \emph{Say it:} \emph{raskátyvat}
  \item \emph{Say it:} \emph{raskátyvat}, once more
  \item \emph{Say it:} a sentence of your own with \emph{raskátyvat}
  \item \emph{Recall:} say the Russian for dough, then the Russian for to knead, then say \emph{raskátyvat} again
\end{itemize}

\begin{tcolorbox}[breakable,colback=teal!4,colframe=teal!35!black,title={Before you move on}]
What does \ru{раскатывать} mean? (\textquotedblleft{}To roll out\textquotedblright{}.) Say it once more.
\end{tcolorbox}
\section[\texorpdfstring{\ru{лепить} (lepít)}{лепить (lepít)}]{\ru{лепить} (lepít) — to shape, to make (by hand); to model}

\label{lesson:RU-C331-lepit}



\begin{quote}
Before the new one: say the Russian for to knead, then the Russian for to roll out.
\end{quote}

\subsection*{You'll want to know: \ru{лепить}}
\textbf{\ru{лепить}} (\emph{lepít}) — \textquotedblleft{}to shape, to make (by hand); to model\textquotedblright{}.

\begin{quote}
\ru{Вся} \ru{семья} \ru{лепила} \ru{пельмени}.
\end{quote}

\begin{quote}
\emph{Vsya semyá lepíla pelméni.}
\end{quote}

\begin{quote}
The whole family was making dumplings.
\end{quote}

Imperfective; present \emph{\ru{леплю}}, \emph{\ru{лепишь}}. Perfective \emph{\ru{слепить}} (\emph{slepít}). \emph{\ru{Лепить} \ru{снеговика}} — to make a snowman.

\begin{cousinweb}
\emph{\ru{лепить}} is inherited from Proto-Slavic \emph{lěpiti} \textquotedblleft{}to smear, to stick\textquotedblright{}, from PIE \emph{leyp-} \textquotedblleft{}to stick\textquotedblright{} (EtymDB).
\end{cousinweb}

\begin{grammarlens}[title={What you've built}]
One more word for telling, explaining and describing at B1.
\end{grammarlens}

\subsection*{Your turn}
\begin{itemize}
\raggedright
  \item \emph{Say it:} \emph{lepít}
  \item \emph{Say it:} \emph{lepít}, once more
  \item \emph{Say it:} a sentence of your own with \emph{lepít}
  \item \emph{Recall:} say the Russian for to knead, then the Russian for to roll out, then say \emph{lepít} again
\end{itemize}

\begin{tcolorbox}[breakable,colback=teal!4,colframe=teal!35!black,title={Before you move on}]
What does \ru{лепить} mean? (\textquotedblleft{}To shape, to make (by hand); to model\textquotedblright{}.) Say it once more.
\end{tcolorbox}
\section[\texorpdfstring{\ru{начинка} (nachínka)}{начинка (nachínka)}]{\ru{начинка} (nachínka) — a filling, a stuffing}

\label{lesson:RU-C331-nachinka}



\begin{quote}
Before the new one: say the Russian for to roll out, then the Russian for to shape.
\end{quote}

\subsection*{You'll want to know: \ru{начинка}}
\textbf{\ru{начинка}} (\emph{nachínka}) — \textquotedblleft{}a filling, a stuffing\textquotedblright{}.

\begin{quote}
\ru{Начинка} \ru{в} \ru{пирогах} \ru{была} \ru{из} \ru{капусты}.
\end{quote}

\begin{quote}
\emph{Nachínka v pirogákh bylá iz kapústy.}
\end{quote}

\begin{quote}
The filling in the pies was cabbage.
\end{quote}

Feminine.

\begin{cousinweb}
\emph{\ru{начинка}} is built on \emph{\ru{начинить}} \textquotedblleft{}to stuff\textquotedblright{}, which is usually traced to \emph{\ru{чин}} \textquotedblleft{}order\textquotedblright{}: to fill in good order. EtymDB lists no parent.
\end{cousinweb}

\begin{grammarlens}[title={What you've built}]
One more word for telling, explaining and describing at B1.
\end{grammarlens}

\subsection*{Your turn}
\begin{itemize}
\raggedright
  \item \emph{Say it:} \emph{nachínka}
  \item \emph{Say it:} \emph{nachínka}, once more
  \item \emph{Say it:} a sentence of your own with \emph{nachínka}
  \item \emph{Recall:} say the Russian for to roll out, then the Russian for to shape, then say \emph{nachínka} again
\end{itemize}

\begin{tcolorbox}[breakable,colback=teal!4,colframe=teal!35!black,title={Before you move on}]
What does \ru{начинка} mean? (\textquotedblleft{}A filling, a stuffing\textquotedblright{}.) Say it once more.
\end{tcolorbox}
